\documentclass[11pt,a4paper]{article}
% Estilo comum dos PDFs individuais da Sprint 1.
\usepackage[utf8]{inputenc}
\usepackage[T1]{fontenc}
\usepackage[brazilian]{babel}
\usepackage[a4paper,margin=2.35cm]{geometry}
\usepackage{xcolor}
\usepackage{hyperref}
\usepackage{enumitem}
\usepackage{listings}
\usepackage{parskip}
\usepackage{needspace}
\definecolor{azul}{RGB}{20,63,112}
\hypersetup{colorlinks=true,linkcolor=azul,urlcolor=azul}
\setlist[itemize]{leftmargin=*,itemsep=3pt}
\lstset{basicstyle=\ttfamily\small,frame=single,breaklines=true,columns=fullflexible}
\sloppy
\newcommand{\codigo}[1]{\texttt{#1}}
\newcommand{\titulo}[2]{\begin{center}{\LARGE\bfseries #1\par}\vspace{4pt}{\large #2\par}\vspace{8pt}\rule{0.8\textwidth}{0.4pt}\par\vspace{5pt}{\small Matheus Moreli Durão Silva -- Matrícula 12621CBS217\par Cibersegurança - Bacharelado - Noturno\par}\end{center}}
\newcommand{\cardfonte}[1]{\par\smallskip\noindent\textit{Referência: }\href{#1}{Trello - enunciado e checklist do card}.}

\begin{document}\titulo{S1\_01 - Engenharia de requisitos}{Recorte de estudo e análise do problema}\enlargethispage{3\baselineskip}
\section*{O que este card pede}
Antes de programar, é preciso entender qual problema será resolvido. Neste trabalho, o problema é organizar ativos de TI (por exemplo, notebook, servidor e roteador) e registrar vulnerabilidades que podem afetá-los. O card pede uma análise de domínio: quem usa o sistema, quais informações existem, quais ações o sistema precisa oferecer e quais limites ele deve respeitar.
\section*{Como pensar a análise}
\textbf{Requisito} é uma necessidade verificável do sistema. Um requisito funcional responde ``o que o programa faz''; por exemplo, cadastrar um ativo. Um requisito não funcional descreve uma qualidade ou restrição; por exemplo, o programa deve salvar os dados em arquivo JSON e apresentar mensagens compreensíveis. O \textbf{domínio} é o assunto do sistema: aqui, inventário e vulnerabilidades de ativos de TI.

Uma entidade é algo sobre o qual guardamos dados. Neste projeto há duas entidades principais: \textbf{ativo} (identificador, nome, tipo, responsável e localização) e \textbf{vulnerabilidade} (descrição, severidade e ativo relacionado). Um caso de uso descreve uma interação completa: ``o responsável cadastra um ativo'', ``o usuário consulta um ativo'', ``o usuário registra uma vulnerabilidade''.
\section*{Aplicação no projeto final}
O arquivo \codigo{app.py} transforma a análise em opções de menu: cadastrar, listar, buscar, atualizar e remover ativos; registrar e visualizar vulnerabilidades; salvar e gerar resumo. Essa relação é importante: cada função do programa existe para atender a uma necessidade identificada antes do código. O JSON representa as entidades persistidas quando o programa é fechado.
\section*{Checklist demonstrado}
\begin{itemize}
\item Engenharia de requisitos e análise do problema: transformar uma necessidade em comportamentos claros.
\item Requisitos funcionais e não funcionais: separar ações do sistema de suas condições de qualidade.
\item Domínio, ativos e vulnerabilidades: identificar os conceitos reais da área de cibersegurança.
\item Entidades, informações, casos de uso e boas práticas: definir dados, usuários e fluxos antes de implementar.
\end{itemize}
\section*{Evidência produzida}
Este recorte é a síntese técnica solicitada pelo exercício ``Introdução ao projeto''. A especificação detalhada continua nos cards S1\_02 e S1\_03, para que os requisitos e seus critérios possam ser conferidos depois.
\cardfonte{https://trello.com/c/2nV6OG9q}
\end{document}
